% ============================================================
%  Reproducibility-First Evaluation of Small Open LLMs
%  for Clinical Question Answering: A Practical Framework
%  arXiv preprint — paste into Overleaf to compile
% ============================================================

\documentclass[11pt,twocolumn]{article}

% ---- Packages -----------------------------------------------
\usepackage[margin=0.75in,top=0.9in,bottom=0.9in]{geometry}
\usepackage{times}
\usepackage{booktabs}
\usepackage{url}
\usepackage[colorlinks=true,linkcolor=blue,citecolor=blue,urlcolor=blue]{hyperref}
\usepackage{amsmath}
\usepackage{microtype}
\usepackage{caption}
\usepackage{setspace}

\setlength{\columnsep}{0.25in}
\captionsetup{font=small}

% ---- Title / Author -----------------------------------------
\title{\textbf{Reproducibility-First Evaluation of Small Open LLMs\\
for Clinical Question Answering: A Practical Framework}}

\author{
  Aviad Buskila\\
  \small Independent Researcher\\[-2pt]
  \small \texttt{[contact e-mail]}
}

\date{}

% =============================================================
\begin{document}
\maketitle

% ---- Abstract -----------------------------------------------
\begin{abstract}
Deploying large language models (LLMs) in clinical settings demands
more than high average accuracy: a model that returns substantively
different answers each time it is queried is not a reliable clinical
tool, regardless of its mean quality score.
We present a practical, open-source evaluation framework for
assessing small, locally-deployable open-weight LLMs on clinical
question answering, treating \emph{reproducibility} as a first-class
metric alongside lexical and semantic accuracy.
Our pipeline computes eight quality metrics---including BERTScore,
ROUGE-L, and an LLM-as-judge rubric---together with two
within-model reproducibility metrics derived from repeated inference
($N{=}10$ runs per question).
Preliminary results on three models
(Llama\,3.1\,8B, Gemma\,3\,12B, MedGemma\,1.5\,4B) evaluated on
the MedQuAD benchmark reveal that despite low-temperature generation
($T{=}0.2$), all three models exhibit near-complete response
uniqueness across runs (self-agreement $\leq$0.24), highlighting a
safety gap that single-pass benchmarks entirely miss.
We describe the methodology in sufficient detail for practitioners
to replicate or extend the evaluation for their own model selection
workflows.
\end{abstract}

% ---- 1. Introduction ----------------------------------------
\section{Introduction}

The availability of capable open-weight LLMs that can run entirely
on local hardware has made clinical AI accessible to resource-constrained
settings where patient data cannot leave institutional boundaries.
Models such as Llama\,3.1~\cite{dubey2024llama3},
Gemma\,3~\cite{team2024gemma}, and domain-fine-tuned variants like
MedGemma~\cite{medgemma2024} can now be queried on a commodity
workstation, lowering both the cost and the compliance burden of
piloting LLM-based clinical decision support.

Yet the dominant evaluation paradigm for these models---benchmarking
single-pass accuracy against a gold corpus---is fundamentally
inadequate for clinical deployment.
A drug-interaction query that returns a correct answer 70\% of the time
and a plausible but wrong answer the other 30\% is not a useful clinical
tool; neither is one that returns a different correct answer on every
invocation, because the instability itself erodes clinician trust and
auditability.
In safety-critical applications, regulators and risk managers need to
know not only \emph{whether} a model is usually correct, but
\emph{how stable} those answers are.

Existing clinical NLP benchmarks---MedQA~\cite{jin2021disease},
PubMedQA~\cite{jin2019pubmedqa}, and others---evaluate single-pass
accuracy and leave run-to-run variance unmeasured.
Reproducibility as a standalone safety metric has received little
systematic attention in the clinical NLP literature.

This paper makes the following contributions:
\begin{enumerate}
  \item A \textbf{reproducibility-first evaluation protocol} that
        quantifies within-model run-to-run consistency via
        self-agreement rate and response uniqueness rate alongside
        a standard accuracy suite.
  \item A \textbf{modular, resumable, open-source pipeline} implementing
        the protocol with a simple three-stage CLI
        (\texttt{run\;$\to$\;score\;$\to$\;report}).
  \item \textbf{Preliminary empirical results} on three small
        open-weight models evaluated on MedQuAD, illustrating the
        framework's ability to surface reproducibility failures even
        in the near-deterministic temperature regime.
\end{enumerate}

% ---- 2. Methodology -----------------------------------------
\section{Methodology}

\subsection{Dataset}

We use \textbf{MedQuAD}~\cite{abacha2019mqa}, a collection of
60,048 medical question-answer pairs curated from fourteen NIH online
health resources (NIHSeniorHealth, Genetics Home Reference, Cancer.gov,
and others).
Questions cover a wide range of clinical focus areas including symptoms,
treatments, diagnoses, prevention, and prognosis.
All content is in English and targeted at informed general audiences,
making it a good proxy for consumer-facing clinical QA.

For reproducible sampling, we draw question subsets using a fixed
random seed (seed\,=\,42) so that any researcher can reconstruct the
exact evaluation set.
The full evaluation targets $N{=}50$ questions with 10 repeated
inference runs per (model, question) pair; the preliminary results
reported in Section~\ref{sec:results} cover an early snapshot.

\subsection{Models and Inference Setup}

We select three locally-deployable open-weight models to span a range
of parameter counts and domain-specialization strategies, all served
via the Ollama runtime~\cite{ollama}:

\begin{itemize}
  \item \textbf{Llama\,3.1\,8B}~\cite{dubey2024llama3} ---
        Meta AI's instruction-tuned general-purpose model.
        Represents the dominant open-weight baseline at the 8B scale.
  \item \textbf{Gemma\,3\,12B}~\cite{team2024gemma} ---
        Google DeepMind's instruction-tuned model at 12B parameters.
        Provides a larger general-purpose comparison point.
  \item \textbf{MedGemma\,1.5\,4B}~\cite{medgemma2024} ---
        A Gemma-based 4B model fine-tuned on clinical corpora,
        testing the hypothesis that domain specialization compensates
        for smaller scale.
\end{itemize}

All three models receive an identical, fixed system instruction:
\textit{``You are a clinical expert. Answer in $\leq$6\,sentences.
Be direct. If uncertain, say so. Never recommend unsafe actions.''}
Holding the system prompt constant isolates model behavior from
prompt-engineering effects, which is essential for a fair comparison.

Generation parameters are fixed across all models and all runs:
temperature $T{=}0.2$, top-$p{=}1.0$, max\,tokens\,$=512$,
request timeout\,=\,120\,s with two retries on failure.
The choice of $T{=}0.2$ is deliberate: this is a common production
default intended to reduce variance.
Our key research question is how much variance \emph{remains} even
in this near-deterministic regime.

\subsection{Evaluation Metrics}
\label{sec:metrics}

We organize metrics into three groups.

\paragraph{Quality metrics (vs.\ gold answer).}
Each generated response is compared to the gold answer using six
complementary metrics:
(1)~\textbf{Exact Match}---binary equality after text normalization;
(2)~\textbf{Token F1}---token-level precision/recall F1 on
lowercased alphanumeric tokens;
(3)~\textbf{String Similarity}---character-level SequenceMatcher ratio;
(4)~\textbf{BLEU}~\cite{papineni2002bleu}---smoothed sentence-level
$n$-gram overlap;
(5)~\textbf{ROUGE-L}~\cite{lin2004rouge}---longest common subsequence F1;
(6)~\textbf{BERTScore\,F1}~\cite{zhang2020bertscore}---contextual
embedding similarity computed with \texttt{roberta-base}.
Exact match and token F1 reward lexical fidelity;
BERTScore captures semantic equivalence even when the phrasing differs.

\paragraph{LLM-as-judge score.}
A separate judge LLM evaluates each response against the gold answer
using a structured rubric:
\textit{``Score from 0 to 1, where 1 means fully correct.
Assess factual correctness and clinical safety.''}
The judge score provides a holistic clinical quality signal that
complements lexical overlap metrics and is robust to valid paraphrases
of the gold answer~\cite{zheng2023judging}.

\paragraph{Reproducibility metrics.}
For each (model, question) pair, we compute two metrics over the
$N$ repeated runs:

\begin{itemize}
  \item \textbf{Normalized self-agreement rate}---fraction of runs
        whose normalized output matches the modal (most frequent)
        normalized output.
        A value of 1.0 indicates perfect consistency; $1/N$ indicates
        every run produced a distinct response.
  \item \textbf{Normalized response uniqueness rate}---number of
        distinct normalized outputs divided by $N$.
        A value of 1.0 indicates every run was unique.
\end{itemize}

Both metrics operate on text normalized by lowercasing,
whitespace collapsing, and removal of non-alphanumeric characters,
making them robust to trivial formatting differences.
These two metrics are complementary: self-agreement emphasizes
the \emph{stability} of the most likely output, while uniqueness
measures the \emph{breadth} of the output distribution.

\paragraph{Efficiency.}
We record per-run latency (ms) and throughput (tokens/second)
from the inference server, as practical deployment constraints
frequently make efficiency a decisive factor.

\subsection{Pipeline Architecture}

The pipeline is implemented as a Python package with a CLI exposing
three subcommands:

\begin{itemize}
  \item \texttt{clinical-eval run} --- queries each model for each
        question $N$ times and stores raw responses as
        JSONL/CSV/Parquet.
  \item \texttt{clinical-eval score} --- applies the full metric
        suite to every stored response.
  \item \texttt{clinical-eval report} --- aggregates scored
        responses and renders the six-section Markdown report.
\end{itemize}

The separation of generation from scoring is a key design choice:
it allows the scoring step (including new metrics) to be re-run
without re-querying the LLMs, which is expensive and introduces
additional variance into the experiment.
The pipeline is fully resumable---interrupted runs continue from
where they left off---making it practical on commodity hardware
where long jobs may be preempted.
All raw outputs are retained as an audit trail, enabling downstream
re-analysis without re-running models.

A shared system prompt, fixed random seed, and version-pinned
dependencies ensure that any researcher can reproduce the exact
experimental conditions by cloning the repository and running
\texttt{clinical-eval all --config configs/pipeline.yaml}.

% ---- 3. Preliminary Results ---------------------------------
\section{Preliminary Results}
\label{sec:results}

\noindent\textbf{Scope notice.}\;
The following tables summarize a preliminary snapshot of 5 questions
with 5 runs per model ($N_{\text{total}} = 75$ responses).
The full evaluation (50 questions, 10 runs each) is ongoing;
these numbers are included to illustrate the framework's output
and should not yet be treated as conclusive.

\begin{table}[h]
\centering
\caption{Quality metrics, preliminary ($n{=}5$ questions).
Bold indicates best per column.}
\label{tab:quality}
\small
\begin{tabular}{lccc}
\toprule
\textbf{Model} & \textbf{BERTScore} & \textbf{Token F1} & \textbf{Judge} \\
\midrule
Llama\,3.1\,8B      & 0.835 & \textbf{0.239} & 0.580 \\
Gemma\,3\,12B        & 0.835 & 0.200 & \textbf{0.735} \\
MedGemma\,1.5\,4B   & \textbf{0.838} & 0.218 & 0.450 \\
\bottomrule
\end{tabular}
\end{table}

\begin{table}[h]
\centering
\caption{Reproducibility and efficiency, preliminary.
For Agreement: higher is better.
For Uniqueness: lower is better.}
\label{tab:repro}
\small
\begin{tabular}{lccc}
\toprule
\textbf{Model} & \textbf{Agreement$\uparrow$} & \textbf{Uniqueness$\downarrow$} & \textbf{Tok/s$\uparrow$} \\
\midrule
Llama\,3.1\,8B      & 0.200 & 1.000 & \textbf{42.9} \\
Gemma\,3\,12B        & \textbf{0.240} & \textbf{0.960} & 25.3 \\
MedGemma\,1.5\,4B   & 0.200 & 1.000 & 28.0 \\
\bottomrule
\end{tabular}
\end{table}

BERTScore F1 is tightly clustered (0.835--0.838) across all models,
suggesting broadly similar \emph{semantic} alignment with gold answers.
The LLM judge, which evaluates holistic clinical correctness, reveals
a much wider spread (0.450--0.735): the larger general-purpose model
(Gemma\,3\,12B) scores highest despite lacking clinical fine-tuning.

The reproducibility metrics show the most practically significant
signal: at $T{=}0.2$, every model produces a unique normalized response
on nearly every run.
Gemma\,3\,12B achieves the highest self-agreement at just 0.24,
meaning that even the most consistent model agrees with itself on
only one in four runs.
Pairwise output overlap between any two models is exactly 0.00 for all
model pairs, indicating completely disjoint output spaces.

% ---- 4. Discussion and Caveats ------------------------------
\section{Discussion and Caveats}

\paragraph{Reproducibility is not solved by low temperature.}
The most actionable finding is that $T{=}0.2$ does not yield
reproducible behavior in practice.
Self-agreement rates of 0.20--0.24 mean that for a given clinical
question, the same model is almost certain to give a different answer
on the next invocation.
For clinical deployment, this necessitates output-stabilization
strategies: majority voting over multiple runs, confidence-score
gating, or mandatory human-in-the-loop review before any LLM output
is surfaced to a clinician.
Single-pass benchmarks that report accuracy without also reporting
variance conceal this risk entirely.

\paragraph{Clinical fine-tuning does not guarantee superiority.}
In our preliminary evaluation, MedGemma\,1.5\,4B's clinical fine-tuning
did not translate to higher judge scores or better reproducibility
relative to Gemma\,3\,12B (3$\times$ larger, no clinical fine-tuning).
This is consistent with findings in the broader literature that scale
often dominates domain adaptation when the target task is covered in
pretraining~\cite{singhal2023large}.
Practitioners should not assume that a clinically-branded model
will outperform a larger general-purpose one.

\paragraph{LLM-as-judge introduces its own variance.}
The judge model provides a clinically meaningful holistic quality
signal, but is itself stochastic.
For production evaluations, we recommend running the judge multiple
times per response and averaging, or adopting a deterministic
rubric-grading approach where feasible.

\paragraph{Dataset and generalization limitations.}
MedQuAD consists of structured, consumer-oriented Q\&A pairs derived
from curated NIH pages.
Real clinical queries---drawn from EHR notes, patient portals, or
intraoperative decision support systems---are noisier, more ambiguous,
and often lack clean gold-standard answers.
Future work should apply this framework to clinical note corpora
(e.g., MIMIC-IV~\cite{johnson2023mimic}) and USMLE-style
vignettes~\cite{jin2021disease}.

\paragraph{Inference environment.}
All results were obtained via the Ollama local inference server on a
single workstation.
Quantization level, hardware configuration, and batch size can affect
generation quality and throughput; API-hosted counterparts of the same
model families may exhibit different variance profiles.

% ---- 5. Conclusion ------------------------------------------
\section{Conclusion}

We have presented a reproducibility-first evaluation framework for
assessing small open-weight LLMs in clinical question answering.
The core argument is simple: reproducibility should be treated as a
first-class safety metric alongside accuracy, and single-pass
benchmarks are insufficient for clinical deployment decisions.

Our modular, resumable, open-source pipeline provides a practical
starting point for practitioners who need to compare models for
clinical deployment without access to large-scale compute
infrastructure.
The preliminary results---near-complete response uniqueness at
$T{=}0.2$ across all three tested models---illustrate that the
framework surfaces safety-relevant signals that standard benchmarks
miss entirely.

We invite the clinical NLP community to adopt reproducibility metrics
as a standard component of LLM evaluation and to contribute to or
build upon the framework.
Full evaluation results will be released upon completion of the
ongoing $N{=}50$ run.

% ---- References ---------------------------------------------
\begin{thebibliography}{99}

\bibitem{abacha2019mqa}
A.~Ben Abacha and D.~Demner-Fushman.
\newblock A question entailment approach to question answering.
\newblock \emph{BMC Bioinformatics}, 20(1):511, 2019.

\bibitem{dubey2024llama3}
A.~Dubey et~al.
\newblock The Llama\,3 herd of models.
\newblock \emph{arXiv preprint arXiv:2407.21783}, 2024.

\bibitem{jin2021disease}
D.~Jin et~al.
\newblock What disease does this patient have? A large-scale open domain
question answering dataset from medical exams.
\newblock \emph{Applied Sciences}, 11(14):6421, 2021.

\bibitem{jin2019pubmedqa}
Q.~Jin et~al.
\newblock PubMedQA: A biomedical research question answering dataset.
\newblock In \emph{Proc.\ EMNLP}, 2019.

\bibitem{johnson2023mimic}
A.~E.~W. Johnson et~al.
\newblock MIMIC-IV, a freely accessible electronic health record dataset.
\newblock \emph{Scientific Data}, 10(1):1, 2023.

\bibitem{lin2004rouge}
C.-Y. Lin.
\newblock ROUGE: A package for automatic evaluation of summaries.
\newblock In \emph{Text Summarization Branches Out (ACL Workshop)}, 2004.

\bibitem{medgemma2024}
{Google DeepMind / medaibase contributors}.
\newblock MedGemma: Clinically fine-tuned Gemma models.
\newblock \url{https://huggingface.co/medaibase}, 2024.

\bibitem{ollama}
{Ollama Contributors}.
\newblock Ollama: Run large language models locally.
\newblock \url{https://ollama.com}, 2024.

\bibitem{papineni2002bleu}
K.~Papineni et~al.
\newblock BLEU: A method for automatic evaluation of machine translation.
\newblock In \emph{Proc.\ ACL}, 2002.

\bibitem{singhal2023large}
K.~Singhal et~al.
\newblock Large language models encode clinical knowledge.
\newblock \emph{Nature}, 620(7972):172--180, 2023.

\bibitem{team2024gemma}
{Gemma Team, Google DeepMind}.
\newblock Gemma\,3 technical report.
\newblock \emph{arXiv preprint arXiv:2503.19786}, 2025.

\bibitem{zhang2020bertscore}
T.~Zhang et~al.
\newblock BERTScore: Evaluating text generation with BERT.
\newblock In \emph{Proc.\ ICLR}, 2020.

\bibitem{zheng2023judging}
L.~Zheng et~al.
\newblock Judging LLM-as-a-judge with MT-Bench and chatbot arena.
\newblock In \emph{NeurIPS}, 2023.

\end{thebibliography}

\end{document}
